\chapter{Referencial teórico}\label{cap:referencial}

Este capítulo revisa as abordagens que explicam por que um megaevento afetaria o emprego, por que o efeito se concentraria na hospitalidade e por que seria temporário, e justifica o arcabouço adotado na dissertação.

\section{Abordagens teóricas possíveis}

\subsection{Impacto econômico de megaeventos: promessas \textit{ex ante} e evidência \textit{ex post}}\label{sub:megaeventos}

Estudos \textit{ex ante}, baseados em matrizes insumo-produto e
multiplicadores, tendem a prever grandes ganhos; Crompton (1995) e Matheson
(2009) mostram as fontes dessa superestimação. No Brasil, Domingues,
Betarelli Junior e Magalhães (2011) projetaram alta de 0,5\% no emprego com a
Copa de 2014. Os estudos \textit{ex post} raramente confirmam esses números.
Hagn e Maennig (2009), comparando mês a mês o desemprego das 12 cidades-sede da
Copa de 2006 com o de 63 cidades alemãs, não encontram efeito significativo, e
observam que a cidade inteira pode ser uma unidade agregada demais para captar
o choque. Mattei e Mello (2017) não encontram diferenças relevantes entre
estados-sede e não sede da Copa de 2014. Os mecanismos recorrentes são a
substituição do gasto local, o \textit{crowding out} de turistas regulares e o
vazamento de renda para fora da região.

\subsection{Efeitos setoriais, temporários e dependentes do contexto}\label{sub:setoriais}

Quando há efeito, ele é setorial e curto. Com dados mensais de 16 setores,
Feddersen e Maennig (2013) encontram aumento do emprego em Atlanta apenas no
mês das Olimpíadas de 1996, concentrado em comércio, alojamento e alimentação e
entretenimento. Boto-García e Santana-Gallego (2025) confirmam, com métodos
recentes, ganhos de emprego em serviços antes das Olimpíadas e Copas, que não
persistem após o evento. O caso mais próximo da COP30 é um evento não
esportivo: Bronzini, Mocetti e Mongardini (2019) mostram que o Jubileu de 2000
elevou a taxa de emprego de Roma, com o ganho concentrado em construção,
comércio e hotéis e restaurantes, sem efeito sobre o valor adicionado per
capita nem sobre o turismo no longo prazo. O contexto também importa: Wan e
Song (2019) encontram resultados distintos para Londres 2012 e para a Copa de
2014 e a Rio 2016, e Kobierecki e Pierzgalski (2022) não identificam efeito
sobre o crescimento das sedes entre 2010 e 2016.

\subsection{Demanda derivada e ajuste do emprego}\label{sub:demanda}

A demanda por trabalho deriva da demanda pelo serviço (HICKS, 1932): um hotel
contrata porque vende diárias. O megaevento é, assim, um choque temporário de
demanda; como as empresas contratam pela demanda esperada, as admissões
antecipam o evento, e os desligamentos seguem a demanda realizada. Alemayehu e
Tveteraas (2020) mostram, com dados diários de hotéis e restaurantes, que o
emprego se ajusta à demanda com defasagem e de forma incompleta, sobretudo em
alimentos e bebidas, o que sugere reações graduais, e não instantâneas, ao fim
do evento.

\subsection{Fluxos de emprego e rotatividade}\label{sub:fluxos}

O emprego resulta de fluxos simultâneos de criação e destruição de postos,
muito maiores que a variação líquida (DAVIS; HALTIWANGER, 1992), que reagem a
choques de demanda (MORTENSEN; PISSARIDES, 1994). Dogru \textit{et al.} (2023)
mostram que a rotatividade na hospitalidade é a mais sensível ao ciclo da
economia. Isso justifica analisar admissões e desligamentos mensais e
controlar pela atividade econômica geral.

\subsection{Legado e economia política urbana}\label{sub:legado}

Preuss (2007) analisa o legado de megaeventos em múltiplas dimensões, e Harvey
(1989) interpreta-os como estratégia de empreendedorismo urbano, com ganhos
concentrados e custos socializados. Essas abordagens enriquecem a
interpretação, mas não geram hipóteses testáveis com dados mensais de emprego.

\section{Escolha do arcabouço teórico}

Adota-se como arcabouço principal a literatura \textit{ex post} sobre
megaeventos (seções~\ref{sub:megaeventos} e~\ref{sub:setoriais}), microfundamentada pela demanda derivada (seção~\ref{sub:demanda})
e pela abordagem de fluxos (seção~\ref{sub:fluxos}); a perspectiva de legado (seção~\ref{sub:legado}) entra na
interpretação. A escolha se justifica pela adequação ao problema (efeito curto,
setorial e medido em fluxos de emprego formal), pela testabilidade (hipóteses
que se traduzem em um modelo com variáveis indicadoras para cada fase) e pelo
descarte fundamentado das alternativas: os estudos \textit{ex ante}
superestimam impactos, e a abordagem crítica opera na escala da governança
urbana. O Quadro~\ref{tab:arcabouco} resume essa correspondência.

\begin{table}[h]
  \centering
  \caption{Escolhas metodológicas e fundamentos teóricos}
  \label{tab:arcabouco}
  \small
  \begin{tabularx}{\textwidth}{Xl}
    \toprule
    \textbf{Escolha metodológica} & \textbf{Fundamento teórico} \\
    \midrule
    Fases: preparação, evento e desmobilização & Megaeventos \textit{ex post} \\
    Foco em Alojamento e Alimentação & Efeitos setoriais (FEDDERSEN; MAENNIG, 2013) \\
    Controle pelo volume de serviços & Demanda derivada; Dogru \textit{et al.} (2023) \\
    Saldo mensal de admissões e desligamentos & Fluxos de emprego e rotatividade \\
    Pará e Região Metropolitana de Belém & Agregação espacial (HAGN; MAENNIG, 2009) \\
    Comparação futura com outras capitais & Contrafactuais (BRONZINI \textit{et al.}, 2019) \\
    \bottomrule
  \end{tabularx}
\end{table}

\textbf{H1}: o saldo de empregos aumenta imediatamente antes e durante o
evento. \textbf{H2}: controlados sazonalidade e volume de serviços, o saldo na
desmobilização fica abaixo do período de referência. \textbf{H3}: os efeitos se
concentram na Região Metropolitana de Belém e no alojamento.

\section{Implicações para a base de dados e para as inferências}

A base usa os microdados do Novo CAGED (admissões e desligamentos formais por
subclasse CNAE e município, pelo mês de competência; BRASIL, 2023) e o índice
de volume de serviços do Pará (IBGE, PMS), de janeiro de 2021 a julho de 2026.
A série começa após a quebra do eSocial em 2020 (IPEA, 2020). Ela permite
inferir se o saldo se desviou do padrão esperado em cada fase e onde o desvio
foi maior, mas, sem grupo de controle, a inferência é de associação
condicional. A literatura aponta o passo seguinte: construir um contrafactual
com outras capitais, por diferenças em diferenças ou controle sintético, como
em Hagn e Maennig (2009) e Bronzini, Mocetti e Mongardini (2019).

